\section{A bundle defined by three graphs}\label{sec:construction}

We construct $G$ together with two line-bundle quotients. Their dual
inclusions will span away from three graphs and agree on those graphs
through specified linear maps. The quotient sequences will prove
ampleness; the graph identities will constrain a quadratic limit on
$G^\vee$.

We choose matrices $A,B$ with two different properties. The relative
lifts $A$, $B$, and $A^{-1}B$ must have finite order, so that their
tensor powers agree at graph intersections. The projective actions of
$A$ and $B$ will have order three and no common fixed point, while
$AB$ will have eigenvalues of distinct moduli. Together these properties
will exclude an invariant measure in Section~\ref{sec:obstruction}. As in the
introduction, $S=\PP^1\times\PP^1$, and we view $\PP^1$ as the lines
in $\CC^2$, with tautological bundle $\OO_{\PP^1}(-1)$.

Put
\begin{equation}\label{eq:matrices}
 a=e^{i\pi/3},\qquad
 s=\frac12-\frac{i(1+2\sqrt2)}{2\sqrt3},\qquad
 A=\begin{pmatrix}a&0\\0&a^{-1}\end{pmatrix},\qquad
 B=\begin{pmatrix}s&1\\ \sqrt2/3&1-s\end{pmatrix}.
\end{equation}
We use the same letters for their induced automorphisms of $\PP^1$.

\begin{lemma}\label{lem:matrix-data}
The matrices $A,B$ belong to $\SL_2(\CC)$ and satisfy
\[
 \tr A=\tr B=1,\qquad
 \tr(A^{-1}B)=-\sqrt2,\qquad
 \tr(AB)=1+\sqrt2.
\]
The orders of $A$, $B$, and $A^{-1}B$ in $\SL_2(\CC)$ are
$6$, $6$, and $8$, respectively. The induced projective orders of $A$ and $B$
are both $3$. The matrix $AB$ has two positive real reciprocal
eigenvalues of distinct moduli, and $A,B$ have no common eigenline.
\end{lemma}

\begin{proof}
The value of $s$ gives
\[
 s(1-s)=1+\frac{\sqrt2}{3},\qquad
 a^{-1}s+a(1-s)=-\sqrt2.
\]
These identities prove the determinant and trace assertions except
for the trace of $AB$. The characteristic polynomial of each of
$A,B$ is $t^2-t+1$, with roots $e^{\pm i\pi/3}$, while that of
$A^{-1}B$ is $t^2+\sqrt2t+1$, with roots $e^{\pm3i\pi/4}$.
The roots in each pair are distinct, proving the claimed orders;
in particular $A^3=B^3=-\Id$.

Cayley--Hamilton gives $A+A^{-1}=\Id$. Multiplication by $B$ and
taking traces therefore yield
\[
 \tr(AB)=\tr B-\tr(A^{-1}B)=1+\sqrt2>2.
\]
Since $\det(AB)=1$, its eigenvalues are positive real reciprocals,
one greater than one. Finally, the eigenlines of $A$ are the two
coordinate lines, and neither is preserved by $B$, whose two
off-diagonal entries are nonzero.
\end{proof}

For $M\in\{I,A,B\}$ let
\[
 \Gamma_M=\{(x,Mx):x\in\PP^1\},\qquad
 D=\Gamma_I\cup\Gamma_A\cup\Gamma_B,
\]
where $D$ has the reduced structure.

\begin{lemma}\label{lem:graph-divisor}
The divisor $D$ has bidegree $(3,3)$ and exactly six singular points,
all ordinary transverse intersections of two components. There are
no triple intersections.
\end{lemma}

\begin{proof}
Each graph is a smooth divisor of bidegree $(1,1)$. An intersection
of $\Gamma_M$ and $\Gamma_N$ corresponds to a fixed point of
$M^{-1}N$, hence to an eigenline of that matrix. The relative
matrices are $A,B,A^{-1}B$ or their inverses. Each has two distinct
eigenvalues, and the multiplier at either projective fixed point is
the ratio of those eigenvalues in one order or the other. This
ratio is not one. The fixed points are therefore simple, so each
pair of graphs meets transversely at two points. A point on all
three graphs would give a common eigenline of $A$ and $B$, contrary
to Lemma~\ref{lem:matrix-data}. Thus the six pairwise intersections
are distinct. Their local equations are products of two regular
parameters, and the asserted description follows.
\end{proof}

Set
\begin{equation}\label{eq:construction-lines}
 b=24,\qquad L_1=p_1^*\OO_{\PP^1}(b),\qquad
 L_2=p_2^*\OO_{\PP^1}(b).
\end{equation}
The exponent $24$ is a common multiple of the orders $6,6,8$ of
the specified relative linear lifts $A,B,A^{-1}B$. It makes the
induced tensor-power maps on the two branches of each graph
intersection agree on their common fiber, as used in the descent
argument below.

Along $\Gamma_M$, the map $M$ on tautological lines induces an
isomorphism from $L_1^{-1}$ to $L_2^{-1}$ by taking its $b$-th
tensor power.

\begin{proposition}\label{prop:construction}
These three isomorphisms glue to an algebraic isomorphism
\begin{equation}\label{eq:alpha}
 \alpha:L_1^{-1}|_D\xrightarrow{\sim}L_2^{-1}|_D,
 \qquad \alpha|_{\Gamma_M}=M^{\otimes b}.
\end{equation}
Let $\beta=(\alpha^\vee)^{-1}:L_1|_D\to L_2|_D$ and define
\begin{equation}\label{eq:equalizer}
 G=\bigl\{(u_1,u_2)\in L_1\oplus L_2:
          u_2|_D=\beta(u_1|_D)\bigr\}.
\end{equation}
Then $G$ is an algebraic vector bundle of rank two. Both projections
$\rho_i:G\to L_i$ are surjective, with exact sequences
\begin{equation}\label{eq:quotient-sequences}
 0\longrightarrow L_i(-D)\longrightarrow G
   \xrightarrow{\rho_j}L_j\longrightarrow0,
 \qquad \{i,j\}=\{1,2\}.
\end{equation}
For $F=G^\vee$, the dual inclusions
$\iota_i=\rho_i^\vee:L_i^{-1}\hookrightarrow F$ span each fiber
over $S\setminus D$ and satisfy
\begin{equation}\label{eq:dual-gluing}
 \iota_2\circ\alpha=\iota_1\qquad\text{on }D.
\end{equation}
\end{proposition}

\begin{proof}
The bundle will be an elementary transformation of $L_1\oplus L_2$
along the effective Cartier divisor $D$, as in
\cite[Introduction, (0.1)]{Maruyama1986}. We verify the gluing and the
local basis directly, since the resulting quotient maps are needed
in both subsequent arguments.

At a point of $\Gamma_M\cap\Gamma_N$, let $v$ span the
first-coordinate line. For some eigenvalue $\lambda$ of $M^{-1}N$
we have $Nv=\lambda Mv$. By Lemma~\ref{lem:matrix-data},
$\lambda^{24}=1$. Consequently
\[
 (Nv)^{\otimes b}=(Mv)^{\otimes b}.
\]
Thus the two maps agree as maps of the actual line fibers. This
also covers the reversed pair, for which the eigenvalue is
$\lambda^{-1}$.

We check descent at a node. If $R$ is the local ring of $S$ there
and $(x),(y)$ are the branch ideals, transversality gives
$(x)\cap(y)=(xy)$ and the exact sequence
\[
 0\longrightarrow R/(xy)\longrightarrow R/(x)\oplus R/(y)
   \longrightarrow R/(x,y)\longrightarrow0,
\]
whose last map is the difference of the residue values. Tensoring
this sequence by the locally free rank-one sheaf
$\operatorname{Hom}(L_1^{-1},L_2^{-1})|_D$ shows that the equal
fiber values are precisely the compatibility needed to descend
the two branch sections. Their common value is nonzero, so the
descended scalar is a unit in the nodal local ring. The component
maps therefore glue to the isomorphism~\eqref{eq:alpha}.

Near a point of $D$, choose frames of $L_1,L_2$ and a local equation
$d$ of $D$. The scalar defining $\beta$ modulo $d$ lifts, after
shrinking, to a regular unit $h$. In these frames the module
\eqref{eq:equalizer} has basis
\begin{equation}\label{eq:local-basis}
 e=(1,h),\qquad f=(0,d).
\end{equation}
Indeed, every element has a unique expression
$u_1e+kf$, since $u_2-hu_1$ is divisible by the non-zero-divisor
$d$. Away from $D$ the sheaf is simply $L_1\oplus L_2$, proving
local freeness everywhere.

The first projection maps $e$ to $1$, and the second maps $e$ to the
unit $h$, so both are surjective. Setting the $j$-th coordinate
equal to zero in~\eqref{eq:equalizer} identifies its kernel with
$L_i(-D)$, proving~\eqref{eq:quotient-sequences}. Their duals are
line subbundles of $F$, and away from $D$ they span. On $D$, if
$v\in L_1^{-1}$ and $(u_1,u_2)\in G$, then
\[
 \bigl(\iota_2(\alpha v)\bigr)(u_1,u_2)
 = (\alpha v)(\beta u_1)=v(u_1)
 =\bigl(\iota_1(v)\bigr)(u_1,u_2),
\]
because $\alpha^\vee\beta=\Id$. This is~\eqref{eq:dual-gluing}.
\end{proof}

For a positive integer $m$, define
\begin{equation}\label{eq:power-maps}
 f_m\bigl([z_0:z_1],[w_0:w_1]\bigr)
   =\bigl([z_0^m:z_1^m],[w_0^m:w_1^m]\bigr).
\end{equation}
This is finite of degree $m^2$, and
$f_m^*\OO_S(r,t)\simeq\OO_S(mr,mt)$. In particular,
$f_m^*P\simeq P^{\otimes m}$. The bundles in the theorem are
\begin{equation}\label{eq:family}
 E_m=f_m^*G\otimes P.
\end{equation}
